\chapter{القياس}\label{ch:g3:measure}

كم الطول، وكم الوزن، وكم يسع الإناء، وكم الساعة؟ القياس يحوّل
العالم إلى أعداد --- بشرط أن يستعمل الجميع الوحدات نفسها. يعرض
هذا الفصل وحدات الطول \omterm{def:g2:measure:units}{والكتلة} \omterm{def:g2:measure:units}{والسعة}، ويعلّم كيف نقرأ
الساعة.

\section{الأطوال}

\begin{definition}[وحدات الطول]\label{def:g3:measure:length}
وحدة الطول الأساسية هي \emph{المتر} (م). وللأطوال الصغيرة:
السنتيمتر (سم) والمليمتر (مم)؛ وللمسافات الطويلة:
الكيلومتر (كم).
\[
1 \text{ م} = 100 \text{ سم}, \qquad
1 \text{ سم} = 10 \text{ مم}, \qquad
1 \text{ كم} = 1\,000 \text{ م}.
\]
\end{definition}

\begin{omfigure}
\begin{tikzpicture}[scale=1.3]
  % ruler body extends past the marks so no numeral is clipped
  \draw[thick] (-0.35,0) rectangle (8.35,0.9);
  \foreach \i in {0,...,8} {
    \draw[thick] (\i,0.9) -- (\i,0.55);
    \node[font=\small] at (\i,0.32) {\i};
  }
  \foreach \i in {0,...,80} \draw (\i*0.1,0.9) -- (\i*0.1,0.72);
  \draw[omThm, very thick] (0,1.25) -- (3.6,1.25);
  \node[omThm, above, font=\small] at (1.8,1.3)
    {$3.6$ سم $= 3$ سم و $6$ مم};
\end{tikzpicture}

{\small قراءة المسطرة: علامات كبيرة للسنتيمترات، وعلامات صغيرة
للمليمترات. والشيء يبدأ عند علامة $0$ --- لا عند حافة
المسطرة! --- وينتهي عند العلامة الصغيرة رقم $6$ بعد $3$.}
\end{omfigure}

\begin{example}[اختيار الوحدة المناسبة]\label{ex:g3:measure:choose}
نملة: نحو $5$ مم. وقلم: نحو $18$ سم. وباب: نحو $2$ م.
ومشي إلى القرية المجاورة: نحو $3$ كم. وقول «طول القلم
$18$» لا يعني شيئًا --- فالوحدة جزء من الجواب.
\end{example}

\begin{method}[تحويل الوحدات]\label{met:g3:measure:convert}
للتحويل، اِضرب أو اِقسم على $10$ أو $100$ أو $1\,000$:
\begin{enumerate}
  \item من وحدة كبيرة إلى صغيرة، \emph{اِضرب}:
        $3$ م $= 3 \times 100 = 300$ سم؛
  \item ومن وحدة صغيرة إلى كبيرة، \emph{اِقسم}:
        $250$ سم $= 250 \div 100 = 2$ م و $50$ سم؛
  \item ولجمع الأطوال، حوّلها إلى الوحدة \emph{نفسها} أولًا:
        $1$ م $+ 35$ سم $= 100 + 35 = 135$ سم.
\end{enumerate}
\end{method}

\section{الكتل والسعات}

\begin{definition}[وحدات الكتلة والسعة]\label{def:g3:measure:mass}
تُقاس الكتل بوحدتَي \emph{الغرام} (غ) و\emph{الكيلوغرام} (كغ)،
وتُقاس السعات (ما يسعه الإناء) بوحدتَي \emph{اللتر} (ل)
و\emph{السنتيلتر} (سل):
\[
1 \text{ كغ} = 1\,000 \text{ غ},
\qquad
1 \text{ ل} = 100 \text{ سل}.
\]
\end{definition}

\begin{example}\label{ex:g3:measure:mass}
رسالة: نحو $20$ غ. وكيس دقيق: $1$ كغ. وحوض استحمام: نحو
$150$ ل. وكأس: نحو $20$ سل. ولموازنة ميزان عليه كيس
$1$ كغ مقابل $650$ غ من التفاح، أضف $1\,000 - 650 = 350$ غ إلى
جهة التفاح.
\end{example}

\section{قراءة الساعة}

\begin{method}[قراءة الوقت]\label{met:g3:measure:clock}
على ساعة ذات عقارب:
\begin{enumerate}
  \item العقرب \emph{القصير} يعطي الساعة: اقرأ العدد الذي
        تجاوزه آخرًا؛
  \item والعقرب \emph{الطويل} يعطي الدقائق: كل علامة كبيرة $5$
        دقائق، فالعقرب الطويل على $3$ يعني $15$ دقيقة؛
  \item وبعد النصف، يمكن أن نقرأ «إلا كذا من الساعة التالية»:
        فالوقت $10{:}50$ هو «الحادية عشرة إلا عشر دقائق».
\end{enumerate}
\end{method}

\begin{omfigure}
\begin{tikzpicture}[scale=1.0]
  % clock 1: 3:00
  \draw[thick] (0,0) circle (1.3);
  \foreach \h in {1,...,12} {
    \node[font=\footnotesize] at ({90-\h*30}:1.05) {\h};
    \draw ({90-\h*30}:1.22) -- ({90-\h*30}:1.3);
  }
  \draw[very thick, omDef] (0,0) -- (0:0.65);
  \draw[thick, omThm] (0,0) -- (90:1.05);
  \node[below, font=\small] at (0,-1.55) {3:00};
  % clock 2: 7:30
  \begin{scope}[xshift=4.4cm]
  \draw[thick] (0,0) circle (1.3);
  \foreach \h in {1,...,12} {
    \node[font=\footnotesize] at ({90-\h*30}:1.05) {\h};
    \draw ({90-\h*30}:1.22) -- ({90-\h*30}:1.3);
  }
  \draw[very thick, omDef] (0,0) -- ({90-7.5*30}:0.65);
  \draw[thick, omThm] (0,0) -- (270:1.05);
  \node[below, font=\small] at (0,-1.55) {7:30};
  \end{scope}
  % clock 3: 10:50
  \begin{scope}[xshift=8.8cm]
  \draw[thick] (0,0) circle (1.3);
  \foreach \h in {1,...,12} {
    \node[font=\footnotesize] at ({90-\h*30}:1.05) {\h};
    \draw ({90-\h*30}:1.22) -- ({90-\h*30}:1.3);
  }
  \draw[very thick, omDef] (0,0) -- ({90-10.83*30}:0.65);
  \draw[thick, omThm] (0,0) -- ({90-10*30}:1.05);
  \node[below, font=\small] at (0,-1.55) {10:50};
  \end{scope}
\end{tikzpicture}

{\small ثلاث ساعات: العقرب الأزرق القصير للساعات، والعقرب
الأحمر الطويل للدقائق. وفي 7:30 يقع عقرب الساعات \emph{بين} 7
و 8؛ وفي 10:50 يكاد يبلغ 11.}
\end{omfigure}

\begin{example}[المدد]\label{ex:g3:measure:duration}
يبدأ الفيلم في $14{:}40$ وينتهي في $16{:}10$. فكم مدته؟
عُدّ في قفزتين:
\begin{enumerate}
  \item من $14{:}40$ إلى $15{:}00$: $20$ دقيقة؛
  \item ومن $15{:}00$ إلى $16{:}10$: $1$ ساعة و $10$ دقائق.
\end{enumerate}
\omterm{def:g1:addition:def}{والمجموع}: $1$ ساعة و $30$ دقيقة. واِنتبه: في الساعة $60$ دقيقة،
فلا يمكن حساب $16{:}10$ ناقص $14{:}40$ كما يُحسب $1610 - 1440$!
\end{example}

\section{تمارين}

\begin{exercise}[$\star$]\label{exo:g3:measure:1}
أي وحدة تستعمل (مم أو سم أو م أو كم) لقياس: عرض مسمار؛
وطول زميل؛ وطول ساحة المدرسة؛ والمسافة بين
مدينتين؟
\end{exercise}

\begin{exercise}[$\star$]\label{exo:g3:measure:2}
حوِّل: $5$ م إلى سم؛ \quad $70$ مم إلى سم؛ \quad $4$ كم إلى م؛
\quad $600$ سم إلى م.
\end{exercise}

\begin{exercise}[$\star$]\label{exo:g3:measure:3}
قِس بمسطرتك (إلى أقرب مليمتر) أضلاع مثلث رسمه زميل لك،
ثم احسب الطول الكلي لحدّه.
\end{exercise}

\begin{exercise}[$\star$]\label{exo:g3:measure:4}
احسب بعد التحويل أولًا: $2$ م $+ 45$ سم (بالسنتيمتر)؛ \quad
$1$ كم $+ 250$ م (بالمتر)؛ \quad $3$ سم $+ 8$ مم (بالمليمتر).
\end{exercise}

\begin{exercise}[$\star$]\label{exo:g3:measure:5}
حوِّل: $3$ كغ إلى غ؛ \quad $2\,500$ غ إلى كغ وغ؛ \quad
$4$ ل إلى سل؛ \quad $350$ سل إلى ل وسل.
\end{exercise}

\begin{exercise}[$\star$]\label{exo:g3:measure:6}
تحتاج وصفة إلى $750$ غ من الدقيق. وعند منى كيس $1$ كغ.
فهل يكفيها؟ وكم يبقى منه؟
\end{exercise}

\begin{exercise}[$\star$]\label{exo:g3:measure:7}
اقرأ الأوقات: العقرب القصير بين $2$ و $3$، والطويل
على $6$؛ والعقرب القصير على $9$، والطويل على $12$؛
والعقرب القصير يكاد يبلغ $5$، والطويل على $10$.
\end{exercise}

\begin{exercise}[$\star$]\label{exo:g3:measure:8}
تبدأ الدراسة في $8{:}30$ وينتهي الصباح في $11{:}45$. فكم يدوم
الصباح؟ (عُدّ بالقفزات، كما في
\cref{ex:g3:measure:duration}.)
\end{exercise}

\begin{exercise}[$\star$]\label{exo:g3:measure:9}
يغادر القطار في $9{:}50$ وتدوم الرحلة $45$ دقيقة. ففي أي
ساعة يصل؟ (اِقفز إلى $10{:}00$ أولًا.)
\end{exercise}

\begin{exercise}[$\star\star$]\label{exo:g3:measure:10}
تسع قنينة $1$ ل. ويملأ سليم كؤوسًا من $20$ سل. فكم كأسًا
كاملة يستطيع أن يصبّ؟ (حوّل أولًا.)
\end{exercise}

\begin{exercise}[$\star\star$]\label{exo:g3:measure:11}
تمشي هناء $600$ م إلى المدرسة، وتعود، كل يوم. فكم مترًا
تمشي في أسبوع دراسي من $5$ أيام؟ أعطِ الجواب بالمتر، ثم
بالكيلومتر.
\end{exercise}
